\scanpage{41}

donnée des commutativités $\alpha_{D^*U}$, et écrire les conditions de compatibilité avec $\alpha,\alpha'$ et $k,k'$. La donnée de $\alpha_{D^*U}$ équivaut à celle de systèmes d'éléments
\[
\boxed{\beta_{i'}\in\mathcal E'}\qquad(i'\in I')
\]
\deleted{\unsure{tels que}}\inserted{tels que} l'on ait, dans les diagrammes
\diagram{\begin{tikzcd}[ampersand replacement=\&]
\unsure{\mathcal G'_{i'}|\widehat G'_{i'}} \arrow[r,"\rho'_{i'}"] \arrow[d,"f_{i'}"'] \& \mathcal E' \arrow[d,"f_{\mathcal E}"] \\
\unsure{\mathcal G_{\tau i'}|\widehat G_{\tau i'}} \arrow[r,"\rho_{\tau i'}"'] \& \mathcal E
\end{tikzcd}}
la relation
\[
f_{\mathcal E}\rho'_{i'}=\operatorname{int}(\beta_{i'})\rho_{\tau i'}\circ f_{i'}.
\]

Reste à exprimer les \inserted{deux} compatibilités \deleted{avec} de
$(\unsure{f_{D^*}},f_D,f_U,\unsure{\alpha_{D^*D}},\alpha_{De},\alpha_{D^*U},\alpha_{Ue})$
\illegible{1} \unsure{ci-dessus} avec $K$ --- la deuxième compatibilité est simplement la compatibilité des $f_{i'}$ avec les $k'_{i'},k_i$, i.e.
\[
f_{i'}\circ k'_{i'}=k_i\qquad(i=\tau i')
\]
et la première \uline{sauf erreur} s'exprime par la ``commutativité''
\[
\deleted{\illegible{1}}\alpha\gamma'_{i'}
 =\inserted{\unsure{\rho(\beta_{i'})|p(\beta_{i'})}}\gamma_i\alpha_{i'}
 \qquad\forall i'\in I'
\]
\[
\boxed{\alpha_{i'}=\gamma_i^{-1}\unsure{\rho|p}(\beta_{i'}^{-1})\alpha\gamma'_{i'}}
\]
\margin{40}

\scanpage{42}
\mspage{19}

En \unsure{résumé|résumant}, les \unsure{homs}, dans un \inserted{syst. de} \deleted{\illegible{1}} diagrammes comm.
\diagram{\begin{tikzcd}[ampersand replacement=\&, column sep=small]
1 \arrow[r] \& T(\bar K) \arrow[r,"k_i"] \arrow[dr] \& \mathcal E_i \arrow[r,"p_i"] \arrow[dr,"\unsure{\rho_i|f_i}"] \& \Gamma \arrow[dr,"\operatorname{int}(\gamma_i)"] \& \\
\& 1 \arrow[r] \& \pi \arrow[r] \& \mathcal E \arrow[r,"p"'] \& \Gamma
\end{tikzcd}\qquad$(i\in I)$}
d'un système analogue, relatif à un ens. d'indices $I'$, est donné par une application
\[
\tau:I'\longrightarrow I
\]
et pour tout $i'\in I'$, posant $i=\tau(i')$, d'un système de flèches verticales $f_{i'}\inserted{{}:\mathcal E'_{i'}\longrightarrow\mathcal E_i}$, et d'une flèche \inserted{verticale} $f_{\mathcal E}:\mathcal E'\longrightarrow\mathcal E$, enfin \deleted{de \illegible{1}} \deleted{$\alpha\in\Gamma$ et} d'un système d'éléments $\beta_{i'}\in\mathcal E$ et d'un $\alpha\in\Gamma$, \deleted{satisfaisant} s'insérant dans le syst. de diagrammes
\diagram{[Prisme formé de quatre lignes, les deuxième et quatrième décalées vers la droite :
\[
\begin{aligned}
1&\longrightarrow T(\bar K)\xrightarrow{k'_{i'}}\mathcal E'_{i'}\xrightarrow{p'_{i'}}\Gamma,\\
1&\longrightarrow\pi'\longrightarrow\mathcal E'\xrightarrow{p'}\Gamma,\\
1&\longrightarrow T(\bar K)\xrightarrow{k_i}\mathcal E_i\xrightarrow{p_i}\Gamma,\\
1&\longrightarrow\pi\longrightarrow\mathcal E\xrightarrow{p}\Gamma.
\end{aligned}
\]
Flèches obliques descendantes de la première à la deuxième ligne :
$T(\bar K)\xrightarrow{\varepsilon'_{i'}}\pi'$,
$\mathcal E'_{i'}\xrightarrow{\rho'_{i'}}\mathcal E'$,
$\Gamma\xrightarrow{\operatorname{int}(\gamma'_{i'})}\Gamma$.
Flèches obliques descendantes de la troisième à la quatrième ligne :
$T(\bar K)\xrightarrow{\varepsilon_i}\pi$,
$\mathcal E_i\xrightarrow{\rho_i}\mathcal E$,
$\Gamma\xrightarrow{\operatorname{int}(\gamma_i)}\Gamma$.
Flèches verticales : entre les deux $T(\bar K)$, une flèche portant $d_{i'}$;
$\pi'\longrightarrow\pi$ sans étiquette lisible;
$\mathcal E'_{i'}\longrightarrow\mathcal E_i$ portant $\unsure{f_{i'}|f^0_{i'}}$;
$\mathcal E'\xrightarrow{f_{\mathcal E}}\mathcal E$;
entre les $\Gamma$ des première et troisième lignes, $\operatorname{int}(\alpha_{i'})$;
entre les $\Gamma$ des deuxième et quatrième lignes, $\operatorname{int}(\alpha)$.
Près de la flèche oblique inférieure $\mathcal E_i\longrightarrow\mathcal E$, annotation surchargée $\unsure{\beta_{i'}}$.]}

\margin{\uline{NB} Comme les $\mathcal E_i\longrightarrow\mathcal E$ sont injectifs, le $f_{i'}$ est connu quand on connait $f_{\mathcal E}$ et $\beta_{i'}$ (l'existence des $f_{i'}$ est donc une \unsure{condition sur les couples} $f_{\mathcal E},\beta_{i'}$ revenant \unsure{à dire} que $\operatorname{int}(\beta_{i'})^{-1}f_{\mathcal E}\rho'_{i'}(\unsure{\mathcal E'_{i'}})$ \unsure{tombe dans} $\unsure{\rho_i(\mathcal E_i)}$).}

\margin{\uline{NB.} \unsure{On} \illegible{1} les $\gamma'_{i'},\gamma_i$ égaux à $1$ (en choisissant l'image de \illegible{3} sur $\operatorname{Spec}(\bar K)$ via l'isom. de $\unsure{\varphi(\rho(D_{i'}^*))\simeq\varphi(\widetilde U)}$ \illegible{3}). i.e. \deleted{\illegible{1}} $\alpha\gamma'_{i'}=p(\beta_{i'})\gamma_i\alpha_{i'}$.}

\unsure{pour} des
\[
\alpha_{i'}=\gamma_i^{-1}p(\beta_{i'}^{-1})\alpha\gamma'_{i'}
\]
i.e. la face \inserted{verticale} postérieure du prisme est commutative (deux conditions, sur deux carrés), la face \inserted{verticale} antérieure \unsure{aussi} (c'est \uline{une} condition, sur le carré de droite, \unsure{l'autre carré}
\margin{41}

\scanpage{43}

comm., s'en déduit par définition de $\pi'\longrightarrow\pi$ comme induit par $f_{\mathcal E}\ldots$), la face verticale gauche du cube de droite étant commutative modulo l'isom. de commutativité $\operatorname{int}(\beta_{i'})$, et la face verticale droite étant commutative (non seulement, par la condition précédente, sur
\[
\unsure{\Gamma'_{i'}}=\operatorname{Im}(\mathcal E'_{i'}\xrightarrow{p'_{i'}}\Gamma),
\]
mais sur $\Gamma$ tout entier) en vertu de la relation plus précise $(*)$.

\figure{[Trait ondulé horizontal séparant les deux parties de la page.]}

Catégorie \unsure{bordélique} \unsure{précise|précisée} (correspondant à un $\Pi_U$ \uline{connexe}) :

Ces objets sont les \deleted{couples \illegible{2}}\inserted{homs de gr. profinis} \deleted{(\unsure{indexés} par un ens. fini $I$) de l'hom}
\[
\mathcal E\xrightarrow{p}\Gamma
\]
(donnant naissance à une suite exacte :
\[
1\longrightarrow\pi\xrightarrow{\unsure{k|\kappa}}\mathcal E\xrightarrow{p}\Gamma
\]
)
et un ens. \inserted{fini} de sous-groupes (indexé par un ens $I$ d'indices)
\[
\mathcal E_i\overset{\rho_i}{\lhook\joinrel\longrightarrow}\mathcal E,
\]
et d'isomorphismes
\[
\mathcal E_i\cap\pi\unsure{\overset{k_i}{\simeq}|\xrightarrow[\sim]{k_i}}T(\bar K)
\]
\margin{42}

\scanpage{44}
\mspage{20}

Un hom d'un système
$(\mathcal E',p',(\mathcal E'_{i'})_{i'\in I'},(k'_{i'})_{i'\in I'})$
dans un syst.
$(\mathcal E,p,(\mathcal E_i)_{i\in I},(k_i)_{i\in I})$
est donné par \unsure{par} un hom
\[
f:\mathcal E'\longrightarrow\mathcal E
\]
\[
\beta_{i'}\in\mathcal E\quad(i'\in I'),\qquad\alpha\in\Gamma,
\]
tels que l'on ait les conditions :

\figure{[Les symboles $f$, $\beta_{i'}$ et $\alpha$ sont entourés.]}

\margin{Les $\beta_{i'}$ \unsure{pas} uniques (si $\beta_{i'}$ convient, aussi $\unsure{\gamma_i}\beta_{i'}$, avec $\unsure{\gamma_i}\in\operatorname{Im}k_i$). Mais $\alpha$ \illegible{2} ?!}

$1^\circ)$\quad $p\circ f_{\mathcal E}=\operatorname{int}(\alpha)\,p'$

$2^\circ)$\quad $\forall i'\in I',\quad\exists i\in I$ (unique!) tel que
\[
\operatorname{int}(\beta_{i'}^{-1})f(\mathcal E'_{i'})\unsure{\subset|=}\mathcal E_i
\]
et un entier $d_{i'}\in\mathbb N^*$ tel que
\[
\operatorname{int}(\beta_{i'}^{-1})f\,k'_{i'}
 =k_i\circ(d_{i'}\cdot\operatorname{id}_{T(\bar K)})
\]

\margin{NB : $d_{i'}$ est aussi unique\ldots}

\margin{[donc $\operatorname{int}(\beta_{i'}^{-1})\overbrace{\unsure{\operatorname{Im}(fk'_{i'})}}^{L'_{i'}}\subset\underbrace{\operatorname{Im}k_i}_{L_i}$ et \unsure{même} \deleted{\illegible{2}} $\operatorname{int}(\beta_{i'}^{-1})\unsure{L'_{i'}}=d_{i'}L_i$.]}

\figure{[Trait horizontal légèrement ondulé, sur toute la largeur du texte.]}

Je m'aperçois qu'il vaut mieux remplacer les $\beta_{i'}$ par les $\beta_{i'}^{-1}$, i.e. prendre l'iso. de commutation plutôt dans le sens
\[
f_{\mathcal E}\rho'_{i'}\xrightarrow{\beta_{i'}}\rho_i f_{i'}
\]
qu'en sens inverse. De plus, conceptuellement le diagramme
\diagram{\begin{tikzcd}[ampersand replacement=\&]
\& \Pi_{D^*}\arrow[dl,"\rho"']\arrow[dr,"\sigma"] \& \\
\Pi_U\arrow[dr,"\varphi"',""{name=L,inner sep=0pt}] \& \& \Pi_D\arrow[dl,"\psi",""{name=R,inner sep=0pt}]\\
\&\Pi_e\&
\arrow[from=L,to=R,bend left=25,"\alpha"]
\end{tikzcd}}
est trop compliqué, il suffit de se donner
\[
\boxed{\Pi_{D^*}\xrightarrow{\rho}\Pi_U\xrightarrow{\varphi}\Pi_e}
\]
et de définir $\Pi_D$ par factorisation canonique
\margin{43}

\scanpage{45}

de l'hom. de groupoïdes $\Pi_{D^*}\longrightarrow\Pi_e$ en un hom qui induit un iso sur les $\pi_0$ et un épi sur les $\pi_1$, suivi d'un hom qui est (épi sur $\pi_0$, et pour cause!, et qui est) \uline{fidèle}. Alors les $1$-morphismes de
\[
\Pi'_{D^*}\xrightarrow{\rho'}\Pi'_U\xrightarrow{\varphi}\Pi_e
\]
dans
\[
\Pi_{D^*}\xrightarrow{\rho}\Pi_U\xrightarrow{\varphi}\Pi_e
\]
sont les \deleted{\illegible{1}}\inserted{\uline{\unsure{quintuplets|quintuples}}} de $2$ foncteurs \inserted{et iso de foncteurs}
$(f_{D^*},f_U,\alpha_{D^*U},\alpha_{Ue},K$
donnant un diagramme avec données de commutation
\diagram{\begin{tikzcd}[ampersand replacement=\&, column sep=large,row sep=large]
\Pi'_{D^*}\arrow[r,"\rho'"]\arrow[d,"f_{D^*}"']\&
\Pi'_U\arrow[r,"\varphi'"]\arrow[d,"f_U"']\&
\Pi_e\arrow[d,equal]\\
\Pi_{D^*}\arrow[r,"\rho"']\&
\Pi_U\arrow[r,"\varphi"']\&\Pi_e
\end{tikzcd}
[Deux flèches courbes, dirigées vers la droite, à l'intérieur des carrés, portent respectivement $\alpha_{D^*U}$ et $\alpha_{Ue}$.]},

\margin{de plus, \deleted{\illegible{1}} $\pi_1(\varphi\rho)=\pi_1(\Pi_{D^*}/\Pi_e)$ peut se définir comme un syst. local \inserted{de groupes \unsure{abéliens}} sur $\Pi_{D^*}$, et on peut alors écrire
$\boxed{\unsure{\bar k}:(\varphi\rho)^*\unsure{T_K|\mathcal T_K}\simeq\pi_1(\varphi\rho)}$.}

\margin{avec \uline{une} condition sur le morphisme
$\pi_1(\varphi'\rho')\unsure{\simeq|\longrightarrow}\unsure{f_{D^*}^{\,*}}(\pi_1(\varphi\rho))$
induit, qui modulo les iso. de Kummer doit être la multiplication par un \unsure{indice de ramification} (qui est une application $\pi_0(\Pi'_{D^*})\longrightarrow\mathbb N^*$). [NB la multiplication par un \uline{\unsure{entier}} \unsure{positif} \illegible{5} constants\ldots]}

si on a deux tels $1$-morphismes
\[
f=(f_{D^*},f_U,\alpha_{D^*U},\alpha_{Ue})
\quad\text{et}\quad
g=(g_{D^*},g_U,\beta_{D^*U},\beta_{Ue})
\]
un \unsure{$0$-morphisme|homomorphisme} de $f$ dans $g$ est formé d'un couple de $\mu_{D^*},\mu_U$ d'isomorphismes de foncteurs
\[
\mu_{D^*}:f_{D^*}\simeq g_{D^*},\qquad
\mu_U:f_U\longrightarrow g_U
\]
compatibles avec $\alpha_{D^*U},\inserted{\beta_{D^*U},}$ et avec $\alpha_{Ue},\inserted{\beta_{Ue},}$ \deleted{et avec}
\margin{44}

\scanpage{46}
\mspage{21}

\deleted{$k,k'$ modulo multiplication par un indice de ramification} (deux conditions de compatibilité sur les deux carrés).

Si on \deleted{suppose} se borne à des $\unsure{\Pi_U,\Pi'_U}$ connexes, on trouve en choisissant comme plus haut un objet \inserted{$\widetilde U$} de $\Pi_U$ et un isom
\[
\boxed{\varphi(\widetilde U)\overset{\lambda}{\simeq}\Omega}
\]
(où $\Omega=\operatorname{Spec}(\bar K)$ est l'objet-repère de $\Pi_e\ldots$), un groupe $\mathcal E$ et un hom
\[
\mathcal E\xrightarrow{p}\Gamma
\]
(qui est changé par autom. int. si on change $\lambda$, et $\mathcal E$ lui-= remplacé par groupe isomorphe, \unsure{d'isom.} défini \unsure{mod.} iso. int., si on change l'objet de référence $\widetilde U$). De \unsure{=}, choisissant un $\widetilde U_i$ dans chaque composante $i\in\pi_0(\Pi_{D^*})$, et des isom
\[
\lambda_i:\rho(\widetilde U_i)\simeq\widetilde U,
\]
on trouve des groupes $\mathcal E_i$ et des homs
\[
\mathcal E_i\xrightarrow{\rho_i}\mathcal E.
\]
Si on change $\lambda_i$, $\rho_i$ est changé par autom intérieur de $\mathcal E$. Si on change \inserted{\illegible{2}} $\widetilde U_i$, $\mathcal E_i$ est remplacé par groupe isom, défini à aut. int. près.
\margin{45}

\scanpage{47}

Considérons les composés
\diagram{\begin{tikzcd}[ampersand replacement=\&]
\mathcal E_i\arrow[r]\arrow[rr,bend left=30,"p_i"]\&\mathcal E\arrow[r,"p"']\&\Gamma
\end{tikzcd}}
\unsure{d'où noyaux}, \deleted{\illegible{1}} $\unsure{K|k}$ est défini par un syst. d'iso
\[
T(\bar K)\overset{k_i}{\simeq}\operatorname{Ker}p_i\qquad(i\in I)
\]
s'insérant dans une famille de suites exactes
\[
1\longrightarrow T(\bar K)\xrightarrow{k_i}\mathcal E_i\xrightarrow{p_i}\Gamma
\]
(NB l'image de $p_i$ est un ss-groupe ouvert)
$k_i$ étant compatible aux actions de $\mathcal E_i$, quand on fait opérer $\mathcal E_i$ sur $T(\bar K)$ par l'action de $\Gamma$ sur $T(\bar K)$, et sur lui-= par autom intérieurs. Introduisant également $\pi=\operatorname{Ker}p$, on trouve donc des homs de suites exactes
\diagram{$(D)$\quad\begin{tikzcd}[ampersand replacement=\&]
1\arrow[r]\&T(\bar K)\arrow[r,"k_i"]\arrow[d,"\unsure{\varepsilon_i}"']\&\mathcal E_i\arrow[r,"p_i"]\arrow[d,"\rho_i"]\&\Gamma\arrow[d,equal]\\
1\arrow[r]\&\pi\arrow[r,"\unsure{k|\kappa}"']\&\mathcal E\arrow[r,"p"']\&\Gamma
\end{tikzcd}\qquad$i\in I$}

\deleted{Considérons}
On peut dire que $\Pi_U$ définit un ``groupe extérieur'' $[\mathcal E]$, et $\Pi_U\longrightarrow\Pi_e$ un hom de groupes extérieurs
\[
[\mathcal E]\longrightarrow[\Gamma],
\]
de = $\Pi_{D^*}$ définit un syst. de ``groupes extérieurs'' $[\mathcal E_i]$, et \deleted{\unsure{du hom}} $\Pi_{D^*}\longrightarrow\Pi_U$
\margin{46}

\scanpage{48}
\mspage{22}

un syst. d'homs extérieurs
\[
[\mathcal E_i]\longrightarrow[\mathcal E],
\]
\deleted{enfin} mais comment en termes de groupes extérieurs exprimer les données de Kummer $k_i$ ?

Revenant au \inserted{système de} diagramme $D$ (relatif à un ens d'indices $I$.) et à un $D'$ analogue (avec ens d'indices $I'$)
\diagram{$D'$\quad\begin{tikzcd}[ampersand replacement=\&]
1\arrow[r]\&T(\bar K)\arrow[r,"k'_{i'}"]\arrow[d,"\varepsilon'_{i'}"']\&\mathcal E'_{i'}\arrow[r,"p'_{i'}"]\arrow[d,"\rho'_{i'}"]\&\Gamma\arrow[d,equal]\\
1\arrow[r]\&\unsure{\pi|\pi'}\arrow[r,"\unsure{k'|\kappa'}"']\&\mathcal E'\arrow[r,"p'"']\&\Gamma
\end{tikzcd}\qquad$i'\in\unsure{I|I'}$}
un $1$-morphisme \inserted{$f$} de $D'$ dans $D$ s'explicite par

$a)$ Un hom
\[
\boxed{\boxed{f_{\mathcal E}:\mathcal E'\longrightarrow\mathcal E}}
\]
\margin{[\uline{\illegible{1}}]}
\margin{décrit le \deleted{\illegible{2}} foncteur $f_U$}
\unsure{tenant du} choix d'un isom
\[
\nu:f(\widetilde U')\simeq\widetilde U
\]
[un autre choix modifie $f_{\mathcal E}$ par un $\operatorname{int}(\beta)$, $\beta\in\mathcal E$)]

$b)$ Une \deleted{hom} application $\tau:I'\longrightarrow I$, et pour tout $i'\in I'$, posant $i=\tau(i')$, un \deleted{application}\inserted{hom de groupes}
\[
\boxed{\mathcal E'_{i'}\xrightarrow{f_{i'}}\mathcal E_i}
\]
\margin{décrit le foncteur $f_{D^*}$}
\margin{47}

\scanpage{49}

[s'explicite en termes du choix d'un iso
\[
f_{D^*}(\widetilde U'_{i'})\overset{\unsure{\nu_{i'}|\gamma_{i'}}}{\simeq}\widetilde U_i,
\]
et modifié par des $\operatorname{int}(\beta_{i'})$, $\beta_{i'}\in\mathcal E_i$, si on change $\unsure{\nu_{i'}|\gamma_{i'}}$]

$c)$ Une donnée de commutation pour
\diagram{\begin{tikzcd}[ampersand replacement=\&]
\mathcal E'\arrow[r,"p'"]\arrow[d,"f_{\mathcal E}"']\&\unsure{\Gamma|\Gamma'}\arrow[d,equal]\\
\mathcal E\arrow[r,"p"']\&\Gamma
\end{tikzcd}\quad
[Flèche courbe intérieure dirigée vers la droite, étiquetée $\alpha$, sur une surcharge $\deleted{\illegible{1}}$.]
\quad$\boxed{\boxed{\deleted{\unsure{\beta}}\alpha\in\Gamma}}$}
\margin{$\alpha$ décrit $\alpha_{Ue}$}
i.e.
\[
\deleted{f_{\mathcal E}p'=\operatorname{int}(\unsure{\beta})p}
\qquad\boxed{pf_{\mathcal E}=\operatorname{int}(\alpha)p'}
\]

$d)$ $\forall i'[\in\unsure{I|I'}]$ une donnée de commutation pour
\diagram{\begin{tikzcd}[ampersand replacement=\&]
\mathcal E'_{i'}\arrow[r,"\rho'_{i'}"]\arrow[d,"f_{i'}"']\&\mathcal E'\arrow[d,"f_{\mathcal E}"]\\
\mathcal E_i\arrow[r,"\rho_i"']\&\mathcal E
\end{tikzcd}\quad
[Flèche courbe intérieure dirigée vers la droite, étiquetée $\alpha_{i'}$.]
\quad$\boxed{\alpha_{i'}\in\mathcal E}\quad(i'\in I')$}
\margin{$\alpha_{i'}$ décrit $\alpha_{D^*U}$}
i.e. :
\[
\boxed{\rho_i f_{i'}=\operatorname{int}(\alpha_{i'})f_{\mathcal E}\rho'_{i'}}
\]

Notons que $c,d$ ensemble définissent une donnée de commutation
\diagram{\begin{tikzcd}[ampersand replacement=\&]
\mathcal E'_{i'}\arrow[r,"p'_{i'}"]\arrow[d,"f_{i'}"']\&\Gamma\arrow[d,equal,"\deleted{\operatorname{int}(\illegible{1})}"]\\
\mathcal E_i\arrow[r,"p_i"']\&\Gamma
\end{tikzcd}\quad
[Flèche courbe intérieure dirigée vers la droite, étiquetée $\operatorname{int}(\beta_{i'})$.]}
avec $\beta_{i'}=p(\alpha_{i'})\alpha$
i.e.
\[
p_i f_{i'}=\operatorname{int}(\beta_{i'})p'_{i'}
\]
\margin{48}

\scanpage{50}
\mspage{23}

i.e. on a commutativité dans
\diagram{\begin{tikzcd}[ampersand replacement=\&]
\mathcal E'_{i'}\arrow[r,"p'_{i'}"]\arrow[d,"f_{i'}"']\&\Gamma\arrow[d,"\operatorname{int}(\beta_{i'})"]\\
\mathcal E_i\arrow[r,"p_i"']\&\Gamma
\end{tikzcd}}
donc $f_{i'}$ induit un hom de suites exactes
\diagram{\begin{tikzcd}[ampersand replacement=\&]
1\arrow[r]\&T(\bar K)\arrow[r,"k'_{i'}"]\arrow[d,"f^0_{i'}"']\&\mathcal E'_{i'}\arrow[r,"\unsure{p'_{i'}|f'_{i'}}"]\arrow[d,"f_{i'}"]\&\Gamma\arrow[d,"\operatorname{int}(\beta_{i'})"]\\
1\arrow[r]\&T(\bar K)\arrow[r,"k_i"']\&\mathcal E_i\arrow[r,"p_i"']\&\Gamma
\end{tikzcd}}
\margin{795}
et il faut exprimer la compatibilité de $f^0_{i'}$ avec l'indice de ramification $d_{i'}\in\mathbb N$, on trouve
\[
\boxed{f^0_{i'}=d_{i'}\chi(\beta_{i'})\operatorname{id}_{T(\bar K)}}
\]
où
\[
\chi:\Gamma\longrightarrow\widehat{\mathbb Z}^{\,*}
\]
est le caractère canonique.

Ceci posé, déterminons, pour un $1$-morphisme
\[
f=(f_{\mathcal E},\inserted{\tau,}(f_{i'}),\alpha,(\unsure{\alpha_{i'}|\alpha_i}))
\]
et un autre
\[
\unsure{f'|g'}=(f'_{\mathcal E},\tau',(f'_{i'}),\alpha',(\alpha'_{i'})),
\]
les $0$-morphismes de l'un dans l'autre.
\margin{49}
